%% ma: L12-B01-0019
%% lop: 12
%% chuong: 1
%% bai: 1
%% dang: 12.1.14
%% loai: TLN
%% muc: VD
%% dap_an: "2,41"
%% nguon: "Ảnh giáo viên gửi 11/10/2026 (Đề kiểm tra giữa HK1 Toán 12, THPT Việt Đức 2026-2027, mã đề 0103), Phần III Câu 1"
%% kien_thuc: "Bài 1 (đơn điệu qua dấu đạo hàm), đạo hàm hàm hợp (lớp 11)"
%% trang_thai: cho-duyet
%% ly_do: "Dạng mới đề xuất, chờ giáo viên duyệt."
\begin{ex}
Cho hàm số $f(x)=x^3-3x$. Đặt $g(x)=f(x^2-2x)$. Biết rằng hàm số $g(x)$ đồng biến trên khoảng $(m;+\infty)$. Tìm giá trị nhỏ nhất của $m$. (Làm tròn kết quả đến chữ số hàng phần trăm.)
\shortans{2,41}
\loigiai{$g'(x)=f'(x^2-2x)(2x-2)$ với $f'(u)=3u^2-3$, nên $g'(x)=3\left[(x^2-2x)^2-1\right](2x-2)=6(x-1)^3(x^2-2x-1)$.

$g'(x)=0\Leftrightarrow x=1$ hoặc $x=1\pm\sqrt2$. Xét dấu: $g'(x)>0$ trên $(1-\sqrt2;1)$ và trên $(1+\sqrt2;+\infty)$; $g'(x)<0$ ở các khoảng còn lại. Vậy $g$ đồng biến trên $(m;+\infty)$ khi và chỉ khi $m\ge1+\sqrt2$, nên giá trị nhỏ nhất của $m$ là $1+\sqrt2\approx2{,}41$.}
\end{ex}
